\chapter{Annexe A : Extraits de code complémentaires}
\label{annexe:code}

\agtlettrine{C}ette annexe rassemble trois extraits de code
représentatifs du projet AGT Infra, choisis pour illustrer concrètement
des mécanismes déjà décrits de façon abstraite dans le corps de ce
rapport : la logique de programmation plutôt que son détail exhaustif.
\minisommaire

\section{Orchestration des phases de déploiement et rollback automatique}

L'extrait suivant, tiré de \texttt{deployment\_engine.py}, montre la
construction de la séquence de phases d'un déploiement ainsi que le
déclenchement du retour arrière automatique lorsque la phase de
vérification sanitaire échoue, comportement déjà détaillé au
chapitre~\ref{chap:realisation}.

\begin{lstlisting}[language=Python]
def _build_steps(log_writer, server, credentials, configuration, env_vars):
    return [
        (DeploymentPhase.SSH_CONNECTION,
         lambda: deployment_phases.run_ssh_connection_phase(log_writer, server, credentials)),
        (DeploymentPhase.PRECHECK,
         lambda: deployment_precheck.run_precheck_phase(log_writer, server, credentials, configuration)),
        (DeploymentPhase.GIT_SYNC,
         lambda: deployment_phases.run_git_sync_phase(log_writer, server, credentials, configuration)),
        (DeploymentPhase.ENV_GENERATION,
         lambda: deployment_phases.run_env_generation_phase(
             log_writer, server, credentials, configuration, env_vars
         )),
        (DeploymentPhase.DOCKER_BUILD,
         lambda: deployment_phases.run_build_and_start_phase(log_writer, server, credentials, configuration)),
        (DeploymentPhase.HEALTHCHECK,
         lambda: deployment_phases.run_healthcheck_phase(log_writer, server, credentials, configuration)),
    ]

# Extrait de la boucle d'execution (_execute_job) :
for phase, step_fn in steps:
    execution.current_phase = phase
    await db.commit()
    result = await step_fn()
    if not result.success:
        await deployment_execution_lifecycle.fail_execution(
            db, execution, job, log_writer, started_at, result.last_error_line
        )
        if phase == DeploymentPhase.HEALTHCHECK and deployment_rollback.should_attempt_rollback(execution):
            await deployment_rollback.run_rollback(db, execution, configuration, server, credentials)
        return
    if phase == DeploymentPhase.GIT_SYNC and result.previous_commit_sha:
        execution.previous_commit_sha = result.previous_commit_sha
\end{lstlisting}

Chaque phase est représentée par une fonction anonyme plutôt que par un
appel direct, afin qu'aucune coroutine ne soit construite avant que la
phase précédente n'ait effectivement réussi. Seul un échec survenant
précisément à la phase de vérification sanitaire déclenche un rollback :
un échec à une autre phase (connexion SSH, synchronisation Git) interrompt
simplement l'exécution, sans tentative de retour en arrière, puisqu'aucun
changement n'aurait encore été appliqué au conteneur en service.

\section{Synchronisation Git idempotente}

L'extrait suivant, tiré de \texttt{deployment\_phases.py}, montre la
phase de synchronisation du dépôt Git. Sa particularité est d'être
rejouable sans effet différent, qu'un dépôt existe déjà sur le serveur
cible ou non, et de capturer le commit précédent pour permettre un
rollback ultérieur.

\begin{lstlisting}[language=Python]
command = (
    f'export GIT_SSH_COMMAND="ssh -i {key_path} -o StrictHostKeyChecking=no"; '
    f'mkdir -p {path}; '
    f'if [ -d {path}/.git ]; then '
    f'cd {path} && echo "PREV_COMMIT:$(git rev-parse HEAD)" '
    f'&& git fetch origin {branch} '
    f'&& git reset --hard origin/{branch} '
    f'&& git clean -fd; '
    f'else '
    f'git clone --branch {branch} --single-branch {clone_url} {path} '
    f'&& echo "PREV_COMMIT:none"; '
    f'fi'
)
ok, output = await run_and_log(log_writer, "GIT_SYNC", server, credentials, command)
if not ok:
    return PhaseResult(success=False, last_error_line=last_nonempty_line(output))

previous_commit_sha = _extract_previous_commit(output)
return PhaseResult(success=True, previous_commit_sha=previous_commit_sha)
\end{lstlisting}

La commande distante teste d'abord si un dépôt existe déjà à l'emplacement
cible. Si c'est le cas, elle imprime le commit courant avant de le
réinitialiser sur la dernière version de la branche suivie ; sinon, elle
clone le dépôt pour la première fois. Le marqueur textuel
\texttt{PREV\_COMMIT} est ensuite extrait de la sortie de la commande,
côté serveur applicatif, pour être conservé sur l'exécution : c'est cette
valeur qui permet à un rollback ultérieur de revenir précisément au
dernier commit connu comme stable.

\section{Résolution d'une permission (refus explicite prioritaire)}

L'extrait suivant, tiré de \texttt{deps.py}, implémente l'algorithme de
résolution d'une permission déjà modélisé par le diagramme de séquence du
chapitre~\ref{chap:etude-technique} : un refus explicite l'emporte
toujours sur une autorisation, qu'elle soit directe ou portée par un
rôle.

\begin{lstlisting}[language=Python]
async def check_permission(
    db: AsyncSession,
    user: User,
    permission_key: str,
    scope_type: ScopeType | None = None,
    resource_id: UUID | None = None,
) -> bool:
    # Etape 1 : SUPER_ADMIN, bypass absolu
    if await _is_super_admin(db, user.id):
        return True

    direct_perms = (await db.execute(
        select(DirectPermission)
        .join(Permission, Permission.id == DirectPermission.permission_id)
        .where(DirectPermission.user_id == user.id, Permission.key == permission_key)
    )).scalars().all()

    # Etape 2 : refus explicite, prioritaire sur tout
    for perm in direct_perms:
        if perm.status != PermissionStatus.DENY:
            continue
        if _scope_matches(perm.scope_type, perm.resource_id, resource_id,
                           scope_type, enforce_expected_scope=False):
            return False

    # Etape 3 : autorisation directe
    for perm in direct_perms:
        if perm.status != PermissionStatus.ALLOW:
            continue
        if _scope_matches(perm.scope_type, perm.resource_id, resource_id,
                           scope_type, enforce_expected_scope=True):
            return True

    # Etape 4 : autorisation portee par un role assigne
    assignments = (await db.execute(
        select(RoleAssignment)
        .join(RolePermission, RolePermission.role_id == RoleAssignment.role_id)
        .join(Permission, Permission.id == RolePermission.permission_id)
        .where(RoleAssignment.user_id == user.id, Permission.key == permission_key)
    )).scalars().all()
    for assignment in assignments:
        if _scope_matches(assignment.scope_type, assignment.resource_id, resource_id,
                           scope_type, enforce_expected_scope=True):
            return True

    # Etape 5 : refus par defaut
    return False
\end{lstlisting}

Les cinq étapes suivent exactement l'ordre déjà présenté au
chapitre~\ref{chap:etude-technique} : seul le rôle immuable court-circuite
entièrement la vérification, un refus explicite est toujours contrôlé
avant toute autorisation, et l'absence de toute correspondance se
résout par un refus par défaut plutôt que par une autorisation implicite.

\section{Annexe B : démonstration vidéo}
% TODO : lien(s) YouTube vers la demonstration video d'AGT Infra, a
% publier par Josue avant integration dans cette section.